\section{Discussion}
\label{sec:discussion}

\para{Why is the cliff below 4 bits?} Our data place the cliff between 4 and 3 bits for models up to 3.8B, which matches the \gptq and \awq papers: 4-bit group-128 is near-lossless there, and 3-bit \rtn breaks down~\cite{frantar2023gptq,lin2024awq}. It contradicts the hypothesis that the cliff sits between INT8 and INT4. One explanation is that with 128-element groups, each group has its own scale and offset, so sixteen levels per group are enough to keep most weights close to their original values, while eight levels are not. The per-channel ablation supports this reading: with a single scale per output channel, the same 16 levels behave roughly like grouped 3.5 bits.

\para{Which tasks break first, and why?} The task ranking follows a simple rule: \emph{tasks whose correctness depends on every token being exact degrade most}. Code passes its unit tests only if the whole function is right, and \gsm needs an exact reasoning chain and final number. \trivia needs the retrieval of rare memorized facts, which are plausibly stored in small weight differences that rounding easily perturbs. Tasks with many acceptable outputs degrade least. \ifeval constraints can be satisfied in many ways, stories are judged holistically, and \squad answers can be copied from the passage. This rule also explains why code, not arithmetic, tops the INT4 ranking: code outputs are longer, and any single error is fatal.

\para{What should practitioners do?} First, group-128 INT6 and INT8 are safe for all six task types, and INT5 is nearly so. Second, INT4 needs task-specific validation, especially for code and for models of 3B and below, and a calibrated quantizer such as \awq recovers most of the \rtn-INT4 loss. Third, \rtn INT3 is unusable below 7B. Fourth, accuracy-level evaluations hide two problems. One is answer churn at 5--6 bits, which matters when users expect consistent answers. The other is silent, fluent errors at 4 bits, which downstream filters will not catch. Teams deploying INT4 models for code or factual assistance should run task-specific regression suites rather than relying on perplexity. Surprisingly, \wikitext perplexity was as good a pooled failure detector as task KL (AUC 0.98), because both explode at 3 bits. Neither caught the subtle INT4 code failures: the perplexity ratio was only 1.12--1.21 at INT4 for models that lost 23--29\% on \humaneval.

\para{How should KL be used?} KL is cheap and label-free, and it reliably flags configurations that are already broken. It does not, however, anticipate a cliff one precision step ahead, because it grows geometrically and crosses any fixed threshold together with accuracy. A single threshold is also unsafe across tasks: code fails at about 0.04 nats, while factual QA is healthy at about 0.4. Task-calibrated thresholds, or the \emph{slope} of KL between adjacent levels, are more promising early-warning signals. We leave both to future work.

\subsection{Limitations}
\label{sec:limitations}

\para{Quantizer coverage.} We use weight-only \rtn, which is pessimistic, plus two \awq checkpoints. Running \gptq, \awq or other calibrated quantizers at every bit width, or quantizing activations and the KV cache, may shift the cliffs by about a bit. Our fake-quantized weights are stored in bf16, which slightly re-rounds the 8-bit grid; this effect is negligible.

\para{Model coverage.} We study seven models, and \phithree only at bf16, 4 and 3 bits. There are only two \llama sizes, one \phithree size, and no models above 8B (a 14B bf16 model did not fit on the shared GPUs). Our size conclusions therefore rest heavily on the 0.5B and 7--8B models.

\para{Statistical resolution.} With 100--400 items per task, per-cell confidence intervals are roughly $\pm$4--8 points, so individual INT4 drops below about 5 points are not resolvable. Critical precisions use point estimates. Greedy decoding gives one deterministic output per configuration, so we report no seed variance, and we did not test sensitivity to prompt format.

\para{Metrics and data.} The creative-writing judge (Qwen3-8B-AWQ) shares a family with several evaluated models and was not validated against human ratings. Judge scores compress quality differences, and \ifeval has a floor of about 0.1--0.4 even for gibberish, so the ``robust'' scores on these tasks partly reflect the metric. The flip analysis shows, however, that \ifeval answers churn at 5 bits as much as \gsm answers. \gsm, \trivia, \squad and \humaneval are public and were likely seen in pre-training. This contamination affects absolute scores more than retention.

\para{KL measurement.} We measure KL on teacher-forced bf16 responses (first 60 items per task, at most 256 tokens) using a top-128 lower bound. KL along the quantized model's own trajectory could differ. Finally, the cliff definitions (0.85 retained; a 0.15 drop) are pre-registered but arbitrary. The location of the steepest drop does not depend on them.

\para{Broader impact.} Quantization makes models cheaper and more accessible. Our results show that it can also introduce fluent, confident errors, \eg wrong facts or subtly broken code, that users cannot easily see. We hope task-resolved evaluation helps deployers avoid shipping such silent regressions.
